\documentclass[11pt]{article}
\usepackage[T1]{fontenc}
\usepackage{lmodern}
\usepackage[a4paper,margin=25mm]{geometry}
\usepackage{amsmath,amssymb}
\usepackage[hidelinks]{hyperref}
\setlength{\emergencystretch}{2em}
\newcommand{\loc}[1]{\par\smallskip\noindent{\small\textit{Source: #1}}\par\smallskip}
\begin{document}
\begin{center}
{\large Probability on Trees and Networks\\Exercise 7.7 and its Author Comment}\\[4pt]
Russell Lyons\quad Yuval Peres\\[2pt]
Author paperback version, 19 August 2026
\end{center}

% BEGIN CONTEXT 7.7
\noindent\textbf{Source and scope (editorial).}
This private reader transcribes the original exercise and its complete printed Comment from the fixed author PDF. The exercise appears both in the chapter body, physical p.\ 266 (printed p.\ 249), and in the collected exercises, physical p.\ 289 (printed p.\ 272); these are one target. The Comment is on physical p.\ 638 (printed p.\ 621). Selected notation and the cited bibliography entry are uncounted context. The separately saved errata dated 19 August 2026 has a Chapter 7 entry for printed p.\ 268, outside these excerpts; no erratum applies here.

\medskip\noindent\textbf{Stochastic domination.}
\loc{Physical pp.\ 264--265 (printed pp.\ 247--248), selected definition.}
Given two percolation measures $\mathbf P$ and $\mathbf P'$ on $G$, we say that $\mathbf P$ \textbf{stochastically dominates} $\mathbf P'$, written $\mathbf P\succcurlyeq\mathbf P'$, if there are random variables $\omega$ and $\omega'$ with laws $\mathbf P$ and $\mathbf P'$, respectively, such that $\omega\geq\omega'$ a.s., which we could also write as $\omega\supseteq\omega'$ a.s.\footnote{See Exercise 4.46(a) and Section 10.2 for more on this concept.}
% END CONTEXT 7.7

% BEGIN PRIMARY STATEMENT 7.7
\subsection*{Exercise 7.7}
\loc{Physical p.\ 289 (printed p.\ 272); also physical p.\ 266 (printed p.\ 249).}
Suppose that $\mathbf P_i$ ($i=1,2,3$) are three percolation measures on $2^A$ such that $\mathbf P_1\succcurlyeq\mathbf P_2$ and $\mathbf P_2\succcurlyeq\mathbf P_3$. Show that there exist random variables $\omega_i$ on a common probability space such that $\omega_i\sim\mathbf P_i$ for all $i$ and $\omega_1\geq\omega_2\geq\omega_3$ a.s.
% END PRIMARY STATEMENT 7.7

% BEGIN PRIMARY PROOF 7.7
\subsection*{Complete printed Comment}
\loc{Physical p.\ 638 (printed p.\ 621), Comment 7.7.}
\noindent\textbf{7.7.} Let $\mu_{\omega_2}$ and $\nu_{\omega_2}$ be the conditional distributions of $\omega_1$ and $\omega_3$ given $\omega_2$. Then let $(\omega_1,\omega_2,\omega_3)$ have the law where $(\omega_1,\omega_3)$ has conditional law $\mu_{\omega_2}\times\nu_{\omega_2}$ given $\omega_2$. See Exercise 18.20 of Billingsley (1995) for a general construction from specified conditional distributions. One can also make a similar construction for finite $A$ and then use Kolmogorov's existence theorem.
% END PRIMARY PROOF 7.7

% BEGIN ATTRIBUTION 7.7
\medskip\noindent\textbf{Referenced bibliography (uncounted).}
\loc{Physical p.\ 669 (printed p.\ 652).}
Billingsley, P. (1995) \emph{Probability and Measure}, 3rd ed. Wiley Series in Probability and Mathematical Statistics. John Wiley, New York. MR: 1324786.

\medskip\noindent\textbf{Material boundary (editorial).}
The author Comment is reproduced in full. Its named existence tool is cited rather than reproduced; the cited Billingsley exercise was not obtained for this reader. The support and marginal checks remain implicit exactly as printed. The existing owner, independent and adjudication records retain H2 and \texttt{countable\_complete\_H2=true}, with the owner's medium-confidence qualification about that external existence step retained. No additional proof steps, hypotheses or mathematical certification are supplied. PDF transcription is the source identity; editable TeX is its derived format. Public redistribution remains on hold.
% END ATTRIBUTION 7.7
\end{document}
